\documentclass[10pt,a4paper]{amsart}
\usepackage{amsmath,amssymb,amsthm,graphicx}
% For commutative diagrams we use Xy-pic
\usepackage[all]{xy}

% We use 2cell for 2-commutative diagrams.
\xyoption{2cell}
% Theorem environments.
%
\theoremstyle{plain}
\newtheorem{theorem}[subsection]{Theorem}
\newtheorem{proposition}[subsection]{Proposition}
\newtheorem{lemma}[subsection]{Lemma}

\theoremstyle{definition}
\newtheorem{definition}[subsection]{Definition}
\newtheorem{example}[subsection]{Example}
\newtheorem{exercise}[subsection]{Exercise}
\newtheorem{situation}[subsection]{Situation}

\theoremstyle{remark}
\newtheorem{remark}[subsection]{Remark}
\newtheorem{remarks}[subsection]{Remarks}

\numberwithin{equation}{subsection}

% Macros
%
\def\lim{\mathop{\mathrm{lim}}\nolimits}
\def\colim{\mathop{\mathrm{colim}}\nolimits}
\def\Spec{\mathop{\mathrm{Spec}}}
\def\Hom{\mathop{\mathrm{Hom}}\nolimits}
\def\Ext{\mathop{\mathrm{Ext}}\nolimits}
\def\SheafHom{\mathop{\mathcal{H}\!\mathit{om}}\nolimits}
\def\SheafExt{\mathop{\mathcal{E}\!\mathit{xt}}\nolimits}
\def\Sch{\mathit{Sch}}
\def\Mor{\mathop{\mathrm{Mor}}\nolimits}
\def\Ob{\mathop{\mathrm{Ob}}\nolimits}
\def\Sh{\mathop{\mathit{Sh}}\nolimits}
\def\NL{\mathop{N\!L}\nolimits}
\def\CH{\mathop{\mathrm{CH}}\nolimits}
\def\proetale{{pro\text{-}\acute{e}tale}}
\def\etale{{\acute{e}tale}}
\def\QCoh{\mathit{QCoh}}
\def\Ker{\mathop{\mathrm{Ker}}}
\def\Im{\mathop{\mathrm{Im}}}
\def\Coker{\mathop{\mathrm{Coker}}}
\def\Coim{\mathop{\mathrm{Coim}}}

% Boxtimes
%
\DeclareMathSymbol{\boxtimes}{\mathbin}{AMSa}{"02}

%
% Macros for moduli stacks/spaces
%
\def\QCohstack{\mathcal{QC}\!\mathit{oh}}
\def\Cohstack{\mathcal{C}\!\mathit{oh}}
\def\Spacesstack{\mathcal{S}\!\mathit{paces}}
\def\Quotfunctor{\mathrm{Quot}}
\def\Hilbfunctor{\mathrm{Hilb}}
\def\Curvesstack{\mathcal{C}\!\mathit{urves}}
\def\Polarizedstack{\mathcal{P}\!\mathit{olarized}}
\def\Complexesstack{\mathcal{C}\!\mathit{omplexes}}
% \Pic is the operator that assigns to X its picard group, usage \Pic(X)
% \Picardstack_{X/B} denotes the Picard stack of X over B
% \Picardfunctor_{X/B} denotes the Picard functor of X over B
\def\Pic{\mathop{\mathrm{Pic}}\nolimits}
\def\Picardstack{\mathcal{P}\!\mathit{ic}}
\def\Picardfunctor{\mathrm{Pic}}
\def\Deformationcategory{\mathcal{D}\!\mathit{ef}}
\usepackage{cleveref}
\crefname{originalreference}{original reference}{original references}
\Crefname{originalreference}{Original reference}{Original references}
\makeatletter
\newlabel{lemma-supported-in-closed-points}{{0F1F}{source}{}{sourcepointer-1}{}}
\newlabel{lemma-supported-in-closed-points@cref}{{[originalreference][0][]0F1F}{[1][0][]source}{}{}{}}
\makeatother
\crefname{originallemma}{original lemma}{original lemmas}
\Crefname{originallemma}{Original lemma}{Original lemmas}
\makeatletter
\newlabel{sites-cohomology-lemma-cohomology-over-U-trivial}{{0BKZ}{source}{}{sourcepointer-2}{}}
\newlabel{sites-cohomology-lemma-cohomology-over-U-trivial@cref}{{[originallemma][0][]0BKZ}{[1][0][]source}{}{}{}}
\makeatother
\title{A complex supported at closed points over a separably closed field vanishes if and only if all its cohomology vanishes}
\author{The Stacks Project Authors}
\date{stacks-project:etale-cohomology.tex:local-snapshot-20261011; GNU FDL 1.2 or later}
\begin{document}
\maketitle
\noindent\textit{Editorial note.} All delivered mathematical text is the author's own, copied from the pinned Stacks Project chapter file etale-cohomology.tex; no statement or proof step was rewritten, reordered or removed.\par\medskip
\noindent\textit{Editorial note.} The counted proof is the author's complete argument: the proof is written out in full: the trivial direction, the bounded below case using the lowest nonvanishing cohomology sheaf, and the general case through the quasi-compact objects and the identification for U = X.\par\medskip
\noindent\textit{Editorial note.} The label lemma-vanishing-closed-points names the counted statement, and the proof environment immediately following it in the source is delivered verbatim as its proof.\par\medskip
\noindent\textit{Editorial note.} The counted statement is carried by the author's own lemma environment in the chapter source.\par\medskip
\noindent\textit{Editorial note.} The statement and proof are copied from the pinned chapter file; the two cited inputs are kept as external pointers with their canonical Stacks tags.\par\medskip
\noindent\textit{Editorial note.} The preamble is the chapter's own native preamble excerpt already used for this source: the Xy-pic setup (native preamble lines 14-18) and the native theorem declarations with the mathematical macros (native preamble lines 154-217).\par\medskip
\noindent\textit{Allowed standard prerequisites (editorial).}\begin{itemize}
\item Complexes supported at closed points. Same chapter, the description of complexes whose stalks vanish outside the closed points, and the resulting triviality of their cohomology over the whole scheme. Supplies the identification H\textasciicircum{}i(X, Q) = H\textasciicircum{}0(X, H\textasciicircum{}i(Q)) used in both cases of the proof.
\item Cohomology over quasi-compact objects is trivial. Cohomology on Sites, the lemma computing the cohomology of a complex satisfying the support condition as its degree zero part. Supplies the general-case statement that H\textasciicircum{}q(U, Q) = H\textasciicircum{}0(U, H\textasciicircum{}q(Q)) for all quasi-compact objects U.
\end{itemize}\medskip
\ifdefined\hypertarget\hypertarget{sourcepointer-1}{}\fi
\noindent\textit{Source reference.} Original reference 0F1F: sources/stacks-etale-cohomology/pinned-original-source.tex, source lines 20333--20333. Prerequisite pointer invoked at proof line(s) 20387; it names a standard fact of the same chapter that the counted argument uses, and every step of the counted proof built on it is retained verbatim in the delivered body.\par\medskip
\ifdefined\hypertarget\hypertarget{sourcepointer-2}{}\fi
\noindent\textit{Source reference.} Original lemma 0BKZ: sources/stacks-chow/tags.txt, source lines 14577--14577. Prerequisite pointer invoked at proof line(s) 20394; it names a standard fact of a sibling chapter that the counted argument uses, and every step of the counted proof built on it is retained verbatim in the delivered body.\par\medskip
\begin{lemma}
\label{lemma-vanishing-closed-points}
Let $K$ be a separably closed field. Let $X$ be a scheme of finite
type over $K$. Let $Q \in D(X_\etale)$. Assume that $Q_{\overline{x}}$
is nonzero only if $x$ is a closed point of $X$. Then
$$
Q = 0 \Leftrightarrow H^i(X, Q) = 0 \text{ for all }i
$$
\end{lemma}

\begin{proof}
The implication from left to right is trivial. Thus we need
to prove the reverse implication.

\medskip\noindent
Assume $Q$ is bounded below; this cases suffices for almost all
applications. If $Q$ is not zero, then we can look at the smallest $i$
such that the cohomology sheaf $H^i(Q)$ is nonzero.
By Lemma \ref{lemma-supported-in-closed-points} we have
$H^i(X, Q) = H^0(X, H^i(Q)) \not = 0$ and we conclude.

\medskip\noindent
General case. Let $\mathcal{B} \subset \Ob(X_\etale)$ be the
quasi-compact objects. By Lemma \ref{lemma-supported-in-closed-points}
the assumptions of Cohomology on Sites, Lemma
\ref{sites-cohomology-lemma-cohomology-over-U-trivial}
are satisfied. We conclude that $H^q(U, Q) = H^0(U, H^q(Q))$
for all $U \in \mathcal{B}$. In particular, this holds for $U = X$.
Thus the conclusion by Lemma \ref{lemma-supported-in-closed-points}
as $Q$ is zero in $D(X_\etale)$ if and only if $H^q(Q)$ is zero
for all $q$.
\end{proof}

\end{document}
